% GENERATED FILE. Edit canonical lessons, then run npm run generate:books.
% canonical-source-hash: fnv1a64:fd9b614fe2347b13
% canonical-lessons: GU-C339-asar, GU-C339-prabhav, GU-C339-nipajvu, GU-C339-paristhiti, GU-C339-sanjog

\chapter{Cause and Effect}
\label{ch:gu-cause-and-effect}
\hlchaptermodality{339}

\begin{chapteropening}
\textbf{By the end of this chapter:} \emph{I can say what caused something and what came of it.}

Complete the last lesson of chapter 339: I can say what caused something and what came of it.
\end{chapteropening}

\section[\texorpdfstring{\gu{અસર} (asar)}{અસર (asar)}]{\gu{અસર} (asar) — an effect}

\label{lesson:GU-C339-asar}



\begin{quote}
Before the new one: say the Gujarati for to pester, then the Gujarati for to hope.
\end{quote}

\subsection*{You'll want to know: \gu{અસર}}
\textbf{\gu{અસર}} — \emph{asar} — \textquotedblleft{}an effect\textquotedblright{}.

\begin{quote}
\gu{દવાની} \gu{અસર} \gu{બે} \gu{કલાક} \gu{પછી} \gu{થઈ}.
\end{quote}

\begin{quote}
\emph{davānī asar be kalāk pachī thaī.}
\end{quote}

\begin{quote}
The medicine took effect two hours later.
\end{quote}

A feminine noun.

\begin{cousinweb}
\textbf{\gu{અસર}} (\emph{asar}) is usually traced to Arabic \emph{athar}, \textquotedblleft{}a trace, a mark\textquotedblright{}, which reached Gujarati through Persian. Gujarati uses it for what one thing does to another.
\end{cousinweb}

\begin{grammarlens}[title={What you've built}]
One more word for giving reasons and opinions.
\end{grammarlens}

\subsection*{Your turn}
\begin{itemize}
\raggedright
  \item \emph{Say it:} \emph{asar}
  \item \emph{Say it:} \emph{asar}, once more
  \item \emph{Say it:} \emph{davānī asar be kalāk pachī thaī.} — The medicine took effect two hours later.
  \item \emph{Recall:} say the Gujarati for to pester, then the Gujarati for to hope, then say \emph{asar} again
\end{itemize}

\begin{tcolorbox}[breakable,colback=teal!4,colframe=teal!35!black,title={Before you move on}]
What does \gu{અસર} mean? (\textquotedblleft{}An effect\textquotedblright{}.) Say it once more.
\end{tcolorbox}
\section[\texorpdfstring{\gu{પ્રભાવ} (prabhāv)}{પ્રભાવ (prabhāv)}]{\gu{પ્રભાવ} (prabhāv) — an influence}

\label{lesson:GU-C339-prabhav}



\begin{quote}
Before the new one: say the Gujarati for to hope, then the Gujarati for an effect.
\end{quote}

\subsection*{You'll want to know: \gu{પ્રભાવ}}
\textbf{\gu{પ્રભાવ}} — \emph{prabhāv} — \textquotedblleft{}an influence\textquotedblright{}.

\begin{quote}
\gu{ગુરુનો} \gu{પ્રભાવ} \gu{વિદ્યાર્થીઓ} \gu{પર} \gu{ઘણો} \gu{છે}.
\end{quote}

\begin{quote}
\emph{guruno prabhāv vidyārthīo par ghaṇo che.}
\end{quote}

\begin{quote}
The teacher has a great deal of influence on the students.
\end{quote}

A masculine noun.

\begin{cousinweb}
\textbf{\gu{પ્રભાવ}} (\emph{prabhāv}) is a tatsama, borrowed from Sanskrit \emph{prabhāva}, \textquotedblleft{}might, power\textquotedblright{} (Monier-Williams). Gujarati uses it for the pull one person or thing has on another.
\end{cousinweb}

\begin{grammarlens}[title={What you've built}]
One more word for giving reasons and opinions.
\end{grammarlens}

\subsection*{Your turn}
\begin{itemize}
\raggedright
  \item \emph{Say it:} \emph{prabhāv}
  \item \emph{Say it:} \emph{prabhāv}, once more
  \item \emph{Say it:} \emph{guruno prabhāv vidyārthīo par ghaṇo che.} — The teacher has a great deal of influence on the students.
  \item \emph{Recall:} say the Gujarati for to hope, then the Gujarati for an effect, then say \emph{prabhāv} again
\end{itemize}

\begin{tcolorbox}[breakable,colback=teal!4,colframe=teal!35!black,title={Before you move on}]
What does \gu{પ્રભાવ} mean? (\textquotedblleft{}An influence\textquotedblright{}.) Say it once more.
\end{tcolorbox}
\section[\texorpdfstring{\gu{નીપજવું} (nīpajvũ)}{નીપજવું (nīpajvũ)}]{\gu{નીપજવું} (nīpajvũ) — to result}

\label{lesson:GU-C339-nipajvu}



\begin{quote}
Before the new one: say the Gujarati for an effect, then the Gujarati for an influence.
\end{quote}

\subsection*{You'll want to know: \gu{નીપજવું}}
\textbf{\gu{નીપજવું}} — \emph{nīpajvũ} — \textquotedblleft{}to result\textquotedblright{}.

\begin{quote}
\gu{ગરમીને} \gu{કારણે} \gu{ઘણા} \gu{રોગ} \gu{નીપજે} \gu{છે}.
\end{quote}

\begin{quote}
\emph{garmīne kāraṇe ghaṇā rog nīpje che.}
\end{quote}

\begin{quote}
Many illnesses come from the heat.
\end{quote}

Past (it …): \gu{તે} \gu{રોગ} \gu{ગરમીથી} \gu{નીપજ્યો}, \emph{te rog garmīthī nīpajyo}.

\begin{cousinweb}
\textbf{\gu{નીપજવું}} (\emph{nīpajvũ}) is inherited from Sanskrit \emph{níṣpadyatē}, \textquotedblleft{}falls out, comes forth\textquotedblright{}; CDIAL lists the Gujarati verb, \textquotedblleft{}to be produced\textquotedblright{} (CDIAL 7511).
\end{cousinweb}

\begin{grammarlens}[title={What you've built}]
One more word for giving reasons and opinions.
\end{grammarlens}

\subsection*{Your turn}
\begin{itemize}
\raggedright
  \item \emph{Say it:} \emph{nīpajvũ}
  \item \emph{Say it:} \emph{nīpajvũ}, once more
  \item \emph{Say it:} \emph{garmīne kāraṇe ghaṇā rog nīpje che.} — Many illnesses come from the heat.
  \item \emph{Recall:} say the Gujarati for an effect, then the Gujarati for an influence, then say \emph{nīpajvũ} again
\end{itemize}

\begin{tcolorbox}[breakable,colback=teal!4,colframe=teal!35!black,title={Before you move on}]
What does \gu{નીપજવું} mean? (\textquotedblleft{}To result\textquotedblright{}.) Say it once more.
\end{tcolorbox}
\section[\texorpdfstring{\gu{પરિસ્થિતિ} (paristhiti)}{પરિસ્થિતિ (paristhiti)}]{\gu{પરિસ્થિતિ} (paristhiti) — a situation}

\label{lesson:GU-C339-paristhiti}



\begin{quote}
Before the new one: say the Gujarati for an influence, then the Gujarati for to result.
\end{quote}

\subsection*{You'll want to know: \gu{પરિસ્થિતિ}}
\textbf{\gu{પરિસ્થિતિ}} — \emph{paristhiti} — \textquotedblleft{}a situation\textquotedblright{}.

\begin{quote}
\gu{પૂર} \gu{પછી} \gu{ગામની} \gu{પરિસ્થિતિ} \gu{ખરાબ} \gu{હતી}.
\end{quote}

\begin{quote}
\emph{pūr pachī gāmnī paristhiti kharāb hatī.}
\end{quote}

\begin{quote}
After the flood the village was in a bad way.
\end{quote}

A feminine noun.

\begin{cousinweb}
\textbf{\gu{પરિસ્થિતિ}} (\emph{paristhiti}) is built from Sanskrit \emph{pari-}, \textquotedblleft{}around\textquotedblright{}, and \emph{sthiti}, \textquotedblleft{}standing, a state\textquotedblright{} (Monier-Williams): the state of things all around.
\end{cousinweb}

\begin{grammarlens}[title={What you've built}]
One more word for giving reasons and opinions.
\end{grammarlens}

\subsection*{Your turn}
\begin{itemize}
\raggedright
  \item \emph{Say it:} \emph{paristhiti}
  \item \emph{Say it:} \emph{paristhiti}, once more
  \item \emph{Say it:} \emph{pūr pachī gāmnī paristhiti kharāb hatī.} — After the flood the village was in a bad way.
  \item \emph{Recall:} say the Gujarati for an influence, then the Gujarati for to result, then say \emph{paristhiti} again
\end{itemize}

\begin{tcolorbox}[breakable,colback=teal!4,colframe=teal!35!black,title={Before you move on}]
What does \gu{પરિસ્થિતિ} mean? (\textquotedblleft{}A situation\textquotedblright{}.) Say it once more.
\end{tcolorbox}
\section[\texorpdfstring{\gu{સંજોગ} (sanjog)}{સંજોગ (sanjog)}]{\gu{સંજોગ} (sanjog) — a circumstance}

\label{lesson:GU-C339-sanjog}



\begin{quote}
Before the new one: say the Gujarati for to result, then the Gujarati for a situation.
\end{quote}

\subsection*{You'll want to know: \gu{સંજોગ}}
\textbf{\gu{સંજોગ}} — \emph{sanjog} — \textquotedblleft{}a circumstance\textquotedblright{}.

\begin{quote}
\gu{આવા} \gu{સંજોગોમાં} \gu{આપણે} \gu{શું} \gu{કરવું} \gu{જોઈએ}?
\end{quote}

\begin{quote}
\emph{āvā sanjogomā̃ āpaṇe śũ karvũ joīe?}
\end{quote}

\begin{quote}
What should we do in circumstances like these?
\end{quote}

A masculine noun.

\begin{cousinweb}
\textbf{\gu{સંજોગ}} (\emph{sanjog}) goes back to Sanskrit \emph{saṃyoga}, \textquotedblleft{}a joining, a connection\textquotedblright{} (Monier-Williams), with \emph{y} become \emph{j} as in many words Gujarati inherited. \textquotedblleft{}A circumstance\textquotedblright{} is its Gujarati sense: the way things happen to come together.
\end{cousinweb}

\begin{grammarlens}[title={What you've built}]
One more word for giving reasons and opinions.
\end{grammarlens}

\subsection*{Your turn}
\begin{itemize}
\raggedright
  \item \emph{Say it:} \emph{sanjog}
  \item \emph{Say it:} \emph{sanjog}, once more
  \item \emph{Say it:} \emph{āvā sanjogomā̃ āpaṇe śũ karvũ joīe?} — What should we do in circumstances like these?
  \item \emph{Recall:} say the Gujarati for to result, then the Gujarati for a situation, then say \emph{sanjog} again
\end{itemize}

\begin{tcolorbox}[breakable,colback=teal!4,colframe=teal!35!black,title={Before you move on}]
What does \gu{સંજોગ} mean? (\textquotedblleft{}A circumstance\textquotedblright{}.) Say it once more.
\end{tcolorbox}
